%% rfam_prewarp_paper.tex
%% "Inverse Geometry Compensation and Turntable Exposure Averaging for
%%  Uniform Densification in Radio Frequency Additive Manufacturing"
%% Matthew McCoy, Georgia Institute of Technology
%% Compiled: pdflatex rfam_prewarp_paper.tex
%%
%% Required packages: amsmath, amssymb, graphicx, booktabs, multirow,
%%   hyperref, xcolor, subcaption, float, siunitx, geometry, microtype,
%%   cleveref, setspace, caption

\documentclass[10pt,twocolumn]{article}

%% ---------------------------------------------------------------
%% Packages
%% ---------------------------------------------------------------
\usepackage[letterpaper,
            top=1.0in, bottom=1.0in,
            left=0.75in, right=0.75in,
            columnsep=0.25in]{geometry}
\usepackage{amsmath,amssymb,amsfonts}
\usepackage{graphicx}
\usepackage{booktabs}
\usepackage{multirow}
\usepackage{subcaption}
\usepackage{float}
\usepackage{siunitx}
\usepackage{xcolor}
\usepackage{microtype}
\usepackage{setspace}
\usepackage[hidelinks,
            colorlinks=true,
            linkcolor=blue!60!black,
            citecolor=green!40!black,
            urlcolor=blue!60!black]{hyperref}
\usepackage[capitalise,noabbrev]{cleveref}
\usepackage{caption}

%% ---------------------------------------------------------------
%% Graphics path
%% ---------------------------------------------------------------
\graphicspath{{outputs_eqs/}}

%% ---------------------------------------------------------------
%% SI units configuration
%% ---------------------------------------------------------------
\sisetup{
  per-mode = symbol,
  detect-weight = true,
  detect-family = true
}

%% ---------------------------------------------------------------
%% Caption style
%% ---------------------------------------------------------------
\captionsetup{font=small, labelfont=bf, skip=4pt}

%% ---------------------------------------------------------------
%% Custom commands
%% ---------------------------------------------------------------
\newcommand{\Qrf}{Q_{\mathrm{rf}}}
\newcommand{\bndstd}{\sigma_{\mathrm{bnd}}}
\newcommand{\Leff}{L_{\mathrm{eff}}}
\newcommand{\rhorel}{\rho_{\mathrm{rel}}}
\newcommand{\Tpc}{T_{\mathrm{pc}}}
\newcommand{\phimelt}{\varphi}
\newcommand{\RMSE}{\mathrm{RMSE}}
\newcommand{\EPE}{\mathrm{EPE}}
\newcommand{\vect}[1]{\boldsymbol{#1}}
\newcommand{\nhat}{\hat{\vect{n}}}
\newcommand{\half}{\tfrac{1}{2}}

%% ---------------------------------------------------------------
%% Title / author
%% ---------------------------------------------------------------
\title{\textbf{Inverse Geometry Compensation and Turntable Exposure
Averaging for Uniform Densification in Radio Frequency Additive
Manufacturing}}

\author{Matthew McCoy\\
\small School of Mechanical Engineering, Georgia Institute of Technology\\
\small Atlanta, Georgia 30332, USA\\
\small \texttt{mmccoy@gatech.edu}}

\date{March 2026}

%% ===============================================================
\begin{document}
\maketitle
%% ===============================================================

%% ---------------------------------------------------------------
%% Abstract
%% ---------------------------------------------------------------
\begin{abstract}
Radio Frequency Additive Manufacturing (RFAM) sinters graphite-doped
Nylon~12 powder at \SI{27.12}{\mega\hertz} by selective Joule
dissipation within the doped region.  Because the electro-quasi-static
(EQS) field is inherently non-uniform, the volumetric heating rate
$\Qrf = \half\sigma|\nabla V|^{2}$ varies by more than an order of
magnitude across the part cross-section: electric-field concentration
at geometric corners produces $\Qrf \approx 37\times$ the spatial
mean, normal-field enhancement at the faces perpendicular to the applied
electric field gives $\approx 11\times$, while the faces parallel to
the field couple only at $\approx 1.6\times$ the mean.  These gradients
drive spatial variation in melt fraction and post-sintering density that
causes both geometric distortion and structural non-uniformity in
finished parts.

This paper presents and evaluates three complementary strategies to
mitigate these coupled problems: (1) an edge placement error
(EPE)-based iterative geometry prewarp algorithm, inspired by inverse
lithography technology (ILT) from semiconductor photomask design;
(2) spike assist features, analogous to sub-resolution assist features
(SRAFs) in optical lithography, placed at the under-heated parallel
faces to concentrate field energy locally; and (3) a time-based
turntable rotation schedule that averages the anisotropic exposure by
rotating the part in the fixed RF field.

All three strategies are evaluated on a library of 17 primitive shapes
and in depth on a non-symmetric H-shape target.  The EPE prewarp
reduces geometric boundary RMSE from \SI{0.68}{\milli\meter} to
\SI{0.22}{\milli\meter} for the H-shape in two iterations.  Four $90^\circ$
turntable rotations over the six-minute exposure cycle reduce boundary
density standard deviation by 34\% on the bare H-shape and by 86\%
when combined with prewarp (from $\bndstd = 0.078$ to $0.011$).  The
combined approach achieves a minimum-to-maximum boundary density ratio
of 0.938, representing a factor-of-six reduction in boundary density
non-uniformity compared to the uncompensated baseline.  Spike assist
features provide only modest improvement ($\leq 14\%$ in $\bndstd$)
at geometrically impractical spike heights, limiting their practical
utility.  The synergy between prewarp and turntable rotation exceeds
their independent contributions by 24 percentage points, owing to a
structural interaction between the geometry correction and the exposure
averaging mechanism.
\end{abstract}

\noindent\textbf{Keywords:} Radio frequency additive manufacturing;
inverse lithography; geometry prewarp; edge placement error; turntable
rotation; Nylon~12 sintering; electro-quasi-static field; densification
uniformity.

%% ---------------------------------------------------------------
\section{Introduction}
\label{sec:introduction}
%% ---------------------------------------------------------------

\subsection{The RFAM Process}

Radio Frequency Additive Manufacturing is a powder-bed sintering
process in which selective densification is achieved by incorporating
a conductive dopant---typically graphite---into a thermoplastic powder,
most commonly Nylon~12 (PA12).  The powder bed is placed between two
planar electrodes driven at the Industrial-Scientific-Medical (ISM)
frequency of \SI{27.12}{\mega\hertz}.  In the undoped regions the
powder is nearly lossless ($\sigma \approx 0$, $\varepsilon_r \approx
2$), so virtually no energy is absorbed.  In the doped regions the
effective conductivity rises to $\sigma \approx \SI{0.04}{\siemens\per\meter}$
and the relative permittivity to $\varepsilon_r \approx 20$, causing
selective Joule heating $\Qrf = \half\sigma|\vect{E}|^{2}$.  When the
local temperature exceeds the Nylon~12 melt onset temperature
($\approx\SI{180}{\celsius}$), the powder coalesces under surface
tension and the local density increases from the poured powder density
($\rhorel \approx 0.55$) toward fully dense material ($\rhorel \to
1.0$).

RFAM occupies a distinct niche among additive sintering processes.
Unlike laser powder bed fusion (LPBF), which scans a focused beam
across the layer surface, RFAM energizes the \emph{entire} doped
region simultaneously.  This enables rapid, potentially scalable
throughput but introduces the fundamental challenge addressed in this
paper: the RF electromagnetic field is spatially non-uniform, and
this non-uniformity is encoded directly into the sintered density
distribution.

\subsection{The Two Coupled Problems}

Two distinct but related failure modes arise from RFAM's non-uniform
field:

\paragraph{Geometric distortion.}
Because densification implies volumetric shrinkage, and because the
degree of densification varies spatially, different parts of the
sintered boundary shrink by different amounts.  A nominally square
$20 \times \SI{20}{\milli\meter}$ target will emerge with corners
over-sintered and face midpoints under-sintered, producing a
systematically distorted outline.  Without compensation, the
root-mean-square boundary error is approximately \SI{0.68}{\milli\meter}
for the H-shape studied here---large enough to matter for any
fit-critical application.

\paragraph{Density non-uniformity.}
Even if geometric distortion is compensated by modifying the input
polygon, the underlying heating non-uniformity remains, producing
a part whose boundary density spans a factor of $\approx 1.46\times$
from the best-sintered corners to the worst-sintered parallel-field
face midpoints.  This density gradient is a structural defect that
cannot be corrected post hoc by shape changes alone.

\subsection{Prior Work}

The RFAM process was introduced and characterized physically by
Allison~\cite{allison2015rfam} and by Seepersad and
colleagues~\cite{seepersad2014volumetric,seepersad2019rfam}, who
established the relationship between graphite doping concentration,
RF power, and sintering quality for Nylon~12.  Computational
simulation of the coupled electromagnetic-thermal-densification
problem was developed by McCoy~\cite{mccoy2025rfam_sim}, demonstrating
that a 2D EQS model with uniform $\sigma$ and a $2\%$ generator
coupling efficiency reproduces the spatial pattern of COMSOL
full-wave simulations at significantly lower computational cost.

The concept of compensating for systematic process-induced distortions
by pre-modifying the design geometry is well established in optical
semiconductor lithography.  Optical proximity correction (OPC) and,
more recently, inverse lithography technology (ILT) iteratively modify
the photomask pattern so that diffraction, aberration, and resist
chemistry combine to print the desired target
feature~\cite{Wong2001ilt,Poonawala2007ilt,Shen2009ilt}.  The edge
placement error (EPE) metric quantifies the deviation between the
predicted printed boundary and the target boundary, and EPE-gradient
descent is the canonical algorithm for ILT mask
optimization~\cite{Jia2014ilt}.  To the authors' knowledge, this
approach has not previously been applied to RF-driven sintering.

Turntable-based exposure averaging in powder-bed sintering has been
discussed qualitatively in the context of microwave
sintering~\cite{Clark1996microwave} but has not been studied
quantitatively or combined with geometric compensation for RFAM.

\subsection{Contributions of This Paper}

This paper makes the following specific contributions:
\begin{enumerate}
\item A validated 2D EQS + thermal + densification forward model for
      RFAM with quantified spatial field distributions (\cref{sec:methodology}).
\item An EPE-based iterative prewarp algorithm adapted from ILT for
      arbitrary polygonal RFAM targets, with systematic $\Leff$
      calibration (\cref{sec:prewarp_algorithm}).
\item A systematic simulation study of 17 primitive shape geometries,
      characterizing the relationship between shape symmetry, boundary
      density non-uniformity, and melt fraction
      (\cref{sec:results_shapes}).
\item Analysis of spike assist features---RFAM analogs of
      optical SRAFs---as a means to redistribute heating to under-sintered
      faces (\cref{sec:results_spikes}).
\item A time-based turntable rotation model and quantitative evaluation
      of multi-rotation schedules (\cref{sec:results_turntable}).
\item Demonstration that combined prewarp + turntable rotation achieves
      86\% reduction in boundary density standard deviation on a
      challenging non-symmetric H-shape target
      (\cref{sec:results_combined}).
\end{enumerate}

%% ---------------------------------------------------------------
\section{RFAM Physics and Field Non-Uniformity}
\label{sec:physics}
%% ---------------------------------------------------------------

\subsection{Joule Heating in an EQS Field}

At the RFAM operating frequency of \SI{27.12}{\mega\hertz}, the free-space
wavelength is $\lambda \approx \SI{11.1}{\meter}$, roughly 500 times
larger than the electrode spacing (\SI{22}{\milli\meter}).  The
displacement current term $\partial\vect{D}/\partial t$ is negligible
relative to the conduction current $\vect{J} = \sigma\vect{E}$ within
the doped region for $\sigma \gg \omega\varepsilon_0\varepsilon_r$
(which holds for $\sigma = \SI{0.04}{\siemens\per\meter}$ at
\SI{27.12}{\mega\hertz}).  The field therefore obeys the
electro-quasi-static (EQS) approximation:
\begin{equation}
  \nabla \cdot (\sigma\,\nabla V) = 0
  \label{eq:eqs_pde}
\end{equation}
with the instantaneous electric field $\vect{E} = -\nabla V$.  The
time-averaged volumetric power deposition is:
\begin{equation}
  \Qrf = \half\,\sigma\,|\nabla V|^{2}
  \label{eq:qrf}
\end{equation}
where the factor of $\half$ accounts for time-averaging the sinusoidal
field $V(t) = V_0\cos(\omega t)$.

\subsection{Physical Origin of Field Non-Uniformity}

The spatial variation of $\Qrf$ across a finite doped region can be
understood from three distinct physical mechanisms, each active at a
different part of the boundary.

\paragraph{Corners (37$\times$ mean).}
At a sharp convex corner between two material regions, the electric
field lines must turn through a larger angle than at a flat face.
For a corner of interior half-angle $\alpha$, the field is
amplified approximately as $r^{\pi/\alpha - 1}$ near the tip in a
lossless dielectric~\cite{Meixner1972singular}.  For a $90^\circ$
corner ($\alpha = \pi/2$), the exponent is $+1$, producing a $1/r$
field divergence as $r \to 0$ that concentrates energy strongly at
the tip.  This is the classical electrostatic tip-effect responsible
for lightning rods and field-emission tips, and it dominates the RFAM
heating pattern at all sharp vertices.

\paragraph{Perpendicular faces ($11\times$ mean).}
At an interface normal to the applied field ($z$-direction), the normal
component of $\vect{D} = \varepsilon_r\varepsilon_0\vect{E}$ is
continuous across the boundary.  Therefore:
\begin{equation}
  \varepsilon_{r,\mathrm{doped}}\,E_{n,\mathrm{in}} =
  \varepsilon_{r,\mathrm{powder}}\,E_{n,\mathrm{out}}
  \label{eq:normal_D_continuity}
\end{equation}
giving $E_{n,\mathrm{in}} = (\varepsilon_{r,\mathrm{powder}} /
\varepsilon_{r,\mathrm{doped}})\,E_{n,\mathrm{out}}$.  For
$\varepsilon_{r,\mathrm{doped}} = 20$ and $\varepsilon_{r,\mathrm{powder}} = 2$,
the interior normal field is \emph{reduced} by a factor of 10 compared
to the exterior, but the $\Qrf \propto |E|^{2}$ evaluated inside the
doped region still shows a substantial enhancement over the center
because the exterior field itself is enhanced by the overall
dielectric contrast.  Numerically, the net effect at the top and
bottom face midpoints is $\approx 11\times$ the spatial mean of
$\Qrf$ across the part.

\paragraph{Parallel faces (1.6$\times$ mean).}
At an interface parallel to the applied field, the tangential component
of $\vect{E}$ is continuous:
\begin{equation}
  E_{t,\mathrm{in}} = E_{t,\mathrm{out}}
  \label{eq:tangential_E_continuity}
\end{equation}
There is no field concentration mechanism; the interior field at the
face midpoint is essentially equal to the far-field value of
$\approx 860/60 = \SI{14.3}{\volt\per\milli\meter}$.  The mild
enhancement of $\approx 1.6\times$ arises from the global redistribution
of field lines around the high-permittivity doped region
(depolarization effect), not from any boundary-normal concentration.

\subsection{Consequences for Sintering}

The $37\times$ ratio between corner and center $\Qrf$ means that
corners reach the melt threshold ($\sim\SI{180}{\celsius}$) far earlier
than the face midpoints and substantially earlier than the part
interior.  At the 6-minute canonical exposure time, corners have
$\rhorel \approx 0.80$ while the parallel-field face midpoints barely
reach $\rhorel \approx 0.55$ (initial poured density).  This spatial
gradient in sintering extent is the root cause of both dimensional
distortion and structural non-uniformity.

%% ---------------------------------------------------------------
\section{Methodology}
\label{sec:methodology}
%% ---------------------------------------------------------------

\subsection{Electro-Quasi-Static Forward Model}
\label{sec:eqs_model}

\subsubsection{Governing Equation and Discretization}

\Cref{eq:eqs_pde} is solved on a 2D Cartesian grid representing the
XZ cross-section (horizontal = $x$, vertical = $z$, electrodes on the
$z$-axis).  The domain is $60 \times \SI{60}{\milli\meter}$ with
$120 \times 120$ uniform cells ($\Delta x = \Delta z \approx
\SI{0.504}{\milli\meter}$).  The part occupies a $20 \times
\SI{20}{\milli\meter}$ square at the center; the surrounding
\SI{20}{\milli\meter}-wide annular region is undoped powder.

The discretization uses a standard finite-difference stencil for
\cref{eq:eqs_pde} on the Cartesian grid.  At material interfaces, the
conductivity is set to the harmonic mean of adjacent cell values:
\begin{equation}
  \sigma_{i+\half} =
  \frac{2\,\sigma_i\,\sigma_{i+1}}{\sigma_i + \sigma_{i+1}}
  \label{eq:harmonic_sigma}
\end{equation}
which ensures continuity of normal current $J_n = \sigma E_n$ across
the interface.  The resulting sparse linear system is assembled in
CSR format and solved via \texttt{scipy.sparse.linalg.spsolve}
(SuperLU backend).

\subsubsection{Boundary Conditions}

The electrode boundary conditions are:
\begin{align}
  V &= 0 \quad \text{at } z = +\SI{30}{\milli\meter} \quad
       \text{(top electrode)} \label{eq:bc_top}\\
  V &= V_0 = \SI{860}{\volt} \quad \text{at } z = -\SI{30}{\milli\meter}
       \quad \text{(bottom electrode)} \label{eq:bc_bot}
\end{align}
matching the COMSOL reference model (Tuned\_Sigma.mph).  The lateral
boundaries ($x = \pm\SI{30}{\milli\meter}$) use Neumann conditions
$\partial V/\partial x = 0$ (zero normal current, representing
insulating chamber walls).

\subsubsection{Power Normalization}

The generator delivers $P_{\mathrm{gen}} = \SI{500}{\watt}$ to the
system.  Due to impedance mismatch and coupling losses between the
generator and the powder-bed load, the effective power absorbed by the
doped region is:
\begin{equation}
  P_{\mathrm{eff}} = \eta\,P_{\mathrm{gen}}\,/\,d_{\mathrm{eff}}
  \label{eq:power_eff}
\end{equation}
where $\eta = 0.020$ (2\% coupling efficiency, calibrated to match
COMSOL steady-state temperatures) and $d_{\mathrm{eff}} =
\SI{0.020}{\meter}$ is the effective part depth in the out-of-plane
direction.  After each EQS solve, the $\Qrf$ map is scaled such that
$\int_\Omega \Qrf\,\mathrm{d}A = P_{\mathrm{eff}}$.  This ensures
energy consistency regardless of the instantaneous part geometry and
conductivity distribution.

The EQS solve is repeated every 20 thermal time steps (\SI{10}{\second})
to capture the effect of temperature-dependent material properties.
In practice, $\sigma$ is held constant (the temperature dependence
of graphite-doped Nylon~12 conductivity is not well characterized at
sintering temperatures), so the EQS solution changes only as the
doped-region geometry evolves with the prewarp polygon in the prewarp
algorithm.

\subsection{Thermal Model}
\label{sec:thermal_model}

The transient heat equation with apparent-heat-capacity treatment of
phase change is:
\begin{equation}
  \rho_{\mathrm{mat}}\,c_p^{\mathrm{eff}}\,\frac{\partial T}{\partial t}
  = \nabla\cdot(k\,\nabla T) + \Qrf - h_c\,(T - T_\infty)\,\delta_{\mathrm{top}}
  \label{eq:heat_eqn}
\end{equation}
where $\delta_{\mathrm{top}}$ is 1 on the top boundary and 0
elsewhere.  The apparent specific heat absorbs the latent heat of
fusion $\mathcal{L}$ over a smoothed phase-change interval $\Delta T$:
\begin{equation}
  c_p^{\mathrm{eff}}(T) = c_p + \mathcal{L}\,
  \frac{1}{\Delta T}\,H'\!\left(\frac{T - \Tpc}{\Delta T}\right)
  \label{eq:apparent_cp}
\end{equation}
where $H'$ is a smoothed Heaviside kernel (raised cosine) and
$\Tpc = \SI{180}{\celsius}$, $\Delta T = \SI{10}{\celsius}$,
$\mathcal{L} = \SI{96700}{\joule\per\kilogram}$ for Nylon~12.

Material properties are blended between the solid and liquid phases
based on the local melt fraction $\phimelt \in [0, 1]$, where
$\phimelt = H((T - \Tpc)/\Delta T)$:
\begin{align}
  k(T) &= (1-\phimelt)\,k_s + \phimelt\,k_l \\
  \rho_{\mathrm{mat}}(T) &= (1-\phimelt)\,\rho_s + \phimelt\,\rho_l
\end{align}
Table~\ref{tab:material_props} summarizes all material constants.
Convection is applied only on the top face with $h_c =
\SI{5}{\watt\per\meter\squared\per\kelvin}$ and $T_\infty =
\SI{25}{\celsius}$; all other boundaries are adiabatic.  The implicit
Crank-Nicolson scheme is used with $\Delta t = \SI{0.5}{\second}$ and
5 thermal substeps per EQS interval.

\begin{table}[tb]
\centering
\caption{Material properties for Nylon~12 and graphite-doped powder
         used in the RFAM forward model.}
\label{tab:material_props}
\begin{tabular}{llSl}
\toprule
Property & Symbol & {Value} & {Unit} \\
\midrule
\multicolumn{4}{l}{\textit{Nylon 12 powder (undoped)}} \\
Thermal conductivity & $k_p$ & 0.197 & \si{\watt\per\meter\per\kelvin} \\
Density & $\rho_p$ & 490 & \si{\kilogram\per\meter\cubed} \\
Specific heat & $c_{p,p}$ & 1072 & \si{\joule\per\kilogram\per\kelvin} \\
\midrule
\multicolumn{4}{l}{\textit{Nylon 12 doped solid}} \\
Thermal conductivity & $k_s$ & 0.10 & \si{\watt\per\meter\per\kelvin} \\
Density & $\rho_s$ & 460 & \si{\kilogram\per\meter\cubed} \\
Specific heat & $c_{p,s}$ & 2500 & \si{\joule\per\kilogram\per\kelvin} \\
Electrical conductivity & $\sigma$ & 0.04 & \si{\siemens\per\meter} \\
Relative permittivity & $\varepsilon_r$ & 20 & {---} \\
\midrule
\multicolumn{4}{l}{\textit{Nylon 12 doped liquid}} \\
Thermal conductivity & $k_l$ & 0.26 & \si{\watt\per\meter\per\kelvin} \\
Density & $\rho_l$ & 1010 & \si{\kilogram\per\meter\cubed} \\
Specific heat & $c_{p,l}$ & 3279 & \si{\joule\per\kilogram\per\kelvin} \\
\midrule
\multicolumn{4}{l}{\textit{Phase change}} \\
Melt temperature & $\Tpc$ & 180 & \si{\celsius} \\
Melt range & $\Delta T$ & 10 & \si{\celsius} \\
Latent heat & $\mathcal{L}$ & 96700 & \si{\joule\per\kilogram} \\
\midrule
\multicolumn{4}{l}{\textit{Generator / coupling}} \\
Generator power & $P_{\mathrm{gen}}$ & 500 & \si{\watt} \\
Coupling efficiency & $\eta$ & 0.020 & {---} \\
Effective depth & $d_{\mathrm{eff}}$ & 0.020 & \si{\meter} \\
\bottomrule
\end{tabular}
\end{table}

\subsection{Densification Model}
\label{sec:densification_model}

The relative density field $\rhorel(x,z,t)$ is evolved by a
two-regime Arrhenius model.  In the solid-state sintering regime, neck
growth is driven by surface diffusion and is limited by the melt
fraction:
\begin{equation}
  \left.\frac{\mathrm{d}\rhorel}{\mathrm{d}t}\right|_{\mathrm{ss}}
  = k_{0,\mathrm{ss}}\,\phimelt^{0.8}\,
    \exp\!\left(\frac{-E_{a,\mathrm{ss}}}{RT}\right)(1 - \rhorel)
  \label{eq:densification_ss}
\end{equation}
In the liquid-state sintering regime (viscous flow driven by surface
tension, analogous to Frenkel's model~\cite{Frenkel1945}):
\begin{equation}
  \left.\frac{\mathrm{d}\rhorel}{\mathrm{d}t}\right|_{\mathrm{liq}}
  = k_{0,\mathrm{liq}}\,\phimelt_{\mathrm{liq}}\,
    \exp\!\left(\frac{-E_{a,\mathrm{liq}}}{RT}\right)(1 - \rhorel)
  \label{eq:densification_liq}
\end{equation}
where $\phimelt_{\mathrm{liq}} = \max(\phimelt - 0.5, 0)/0.5$ is the
liquid-phase fraction above the solidus.  The total rate is
$\mathrm{d}\rhorel/\mathrm{d}t = \mathrm{d}\rhorel/\mathrm{d}t|_{\mathrm{ss}}
+ \mathrm{d}\rhorel/\mathrm{d}t|_{\mathrm{liq}}$.

Kinetic parameters are:
$k_{0,\mathrm{ss}} = \SI{5e-3}{\per\second}$,
$E_{a,\mathrm{ss}} = \SI{48000}{\joule\per\mol}$;
$k_{0,\mathrm{liq}} = \SI{8e-2}{\per\second}$,
$E_{a,\mathrm{liq}} = 0$ (viscous flow has negligible Arrhenius
temperature dependence over the narrow sintering temperature window).
The initial density is $\rhorel(t=0) = 0.55$ uniformly throughout
the doped region.

\subsection{EPE-Based Geometry Prewarp Algorithm}
\label{sec:prewarp_algorithm}

\subsubsection{Analogy to Inverse Lithography Technology}

In optical semiconductor lithography, the target pattern on the wafer
is related to the mask pattern by a forward operator $\mathcal{F}$
representing diffraction, lens aberrations, and resist chemistry.  ILT
seeks the mask $M^*$ such that $\mathcal{F}(M^*) \approx T$, where
$T$ is the target wafer pattern.  The standard iterative solution
uses gradient descent on the edge placement error:
$\EPE[i] = (\vect{T}[i] - \hat{\vect{B}}[i]) \cdot \nhat_{\mathrm{out}}[i]$,
where $\hat{\vect{B}}$ is the predicted printed boundary and
$\vect{T}[i]$ is the nearest target boundary point.

For RFAM, the analogous forward operator maps the input polygon $P$
through (a) EQS+thermal sintering simulation, (b) local shrinkage, and
(c) final boundary prediction to yield an estimated sintered boundary
$\hat{\vect{B}}$.  The prewarp algorithm iteratively adjusts $P$ to
minimize $\EPE$ in the same spirit as optical ILT.

\subsubsection{Algorithm}

Let $T$ be the target polygon (desired sintered shape) and $P_k$ be
the prewarp polygon at iteration $k$, both represented as ordered
lists of $N = 512$ equally-spaced vertices.  The algorithm proceeds
as follows:

\begin{enumerate}
\item \textbf{Initialize:} $P_0 = T$.
\item \textbf{Forward simulation:} Run the EQS+thermal+densification
      model on polygon $P_k$ for $t = \SI{360}{\second}$.
\item \textbf{Boundary density sampling:} For each vertex $i$, compute
      the outward normal $\nhat_i$.  Sample $\rhorel$ at a point
      inset $\approx 1.5$ grid cells inward along $\nhat_i$:
      $\rhorel^{\mathrm{bnd}}[i]$.
\item \textbf{Shrinkage prediction:} Estimate the local linear
      shrinkage ratio from the measured boundary density, assuming
      isotropic densification from initial density $\rhorel^{(0)}$:
      \begin{equation}
        s[i] = 1 - \sqrt{\rhorel^{(0)}\,/\,\rhorel^{\mathrm{bnd}}[i]}
        \label{eq:shrinkage_ratio}
      \end{equation}
\item \textbf{Predicted boundary:} Project each vertex outward by the
      shrinkage distance:
      \begin{equation}
        \hat{\vect{B}}[i] = P_k[i] - s[i]\,\Leff\,\nhat_i
        \label{eq:predicted_boundary}
      \end{equation}
      where $\Leff$ is the effective shrinkage length scale
      (calibrated parameter, \cref{sec:leff_calibration}).
\item \textbf{Edge placement error:}
      \begin{equation}
        \EPE[i] = (T[i] - \hat{\vect{B}}[i]) \cdot \nhat_i
        \label{eq:epe}
      \end{equation}
\item \textbf{Polygon update:}
      \begin{equation}
        P_{k+1}[i] = P_k[i] + \alpha\,\EPE[i]\,\nhat_i
        \label{eq:polygon_update}
      \end{equation}
      where $\alpha = 1.0$ is the step size (unit EPE-to-correction
      ratio, justified by the approximately linear relationship between
      input polygon boundary position and sintered boundary position
      for small perturbations).
\item \textbf{Laplacian smoothing:}
      \begin{equation}
        P_{k+1}[i] \leftarrow P_{k+1}[i] +
        \lambda_s\,\bigl(P_{k+1}[i-1] - 2P_{k+1}[i] + P_{k+1}[i+1]\bigr)
        \label{eq:laplacian_smooth}
      \end{equation}
      with $\lambda_s = 0.15$, applied 3 times to suppress
      high-frequency oscillations that arise at vertices where
      $\EPE$ is noisy.
\item \textbf{Resample:} Re-parameterize $P_{k+1}$ to $N = 512$
      equally-spaced vertices by arc-length.
\item \textbf{Convergence check:} Stop if
      $\mathrm{RMS}(\EPE) < \varepsilon_{\EPE}$, with
      $\varepsilon_{\EPE} = \SI{0.25}{\milli\meter}$ for square/circle
      and \SI{0.30}{\milli\meter} for the H-shape.
\end{enumerate}

\subsubsection{$\Leff$ Calibration}
\label{sec:leff_calibration}

The effective shrinkage length $\Leff$ is the only free parameter of
the prewarp algorithm.  Physically, it represents the characteristic
distance over which the boundary density $\rhorel^{\mathrm{bnd}}$
measured at the vertex inset point translates into a linear shrinkage
at the sintered boundary.  It depends on the exposure time, material
properties, and the local $\Qrf$ gradient near the boundary, and
therefore differs between shapes.

Calibration is performed by running the full prewarp algorithm to
convergence for each candidate $\Leff \in \{3, 5, 7, 10, 13, 16\}
\si{\milli\meter}$ and selecting the value that achieves the lowest
post-convergence geometric RMSE:
\begin{equation}
  \RMSE = \sqrt{\frac{1}{N}\sum_{i=1}^{N}\EPE[i]^{2}}
  \label{eq:rmse}
\end{equation}

\subsection{Boundary Density Uniformity Metrics}
\label{sec:uniformity_metrics}

Two complementary scalar metrics quantify the spatial uniformity of
sintering across the part boundary.

\paragraph{Boundary standard deviation.}
The one-cell-thick boundary mask $\mathcal{B}$ is defined as the
set-difference between the part mask and its binary erosion:
\begin{equation}
  \mathcal{B} = \{(i,j) : M_{ij} = 1\} \setminus
                \mathrm{erode}(\{(i,j) : M_{ij} = 1\})
  \label{eq:boundary_mask}
\end{equation}
The boundary standard deviation is then:
\begin{equation}
  \bndstd = \mathrm{std}\bigl(\{\rhorel[i,j] : (i,j)\in\mathcal{B}\}\bigr)
  \label{eq:bnd_std}
\end{equation}
Lower $\bndstd$ indicates more uniform surface sintering.

\paragraph{Min/max ratio.}
\begin{equation}
  r_{\mathrm{bnd}} = \frac{\min_{(i,j)\in\mathcal{B}}\rhorel[i,j]}
                         {\max_{(i,j)\in\mathcal{B}}\rhorel[i,j]}
  \in [0, 1]
  \label{eq:minmax_ratio}
\end{equation}
$r_{\mathrm{bnd}} \to 1$ implies perfectly uniform boundary density.
This metric is complementary to $\bndstd$ in that it captures the
worst-case nonuniformity rather than the distribution spread.

\subsection{Spike Assist Features}
\label{sec:spike_method}

Sub-resolution assist features (SRAFs) in optical lithography are mask
elements that are too small to print but that improve the aerial image
intensity distribution for nearby features they are designed to
assist~\cite{Liebmann2001sraf}.  The RFAM analog is a narrow convex
protrusion at a region of under-heating (a parallel-to-$\vect{E}$ face
midpoint) that concentrates $\vect{E}$ locally through the same tip
effect responsible for corner overheating, with the intent of
increasing $\Qrf$ at the face midpoint without contributing significant
sintered material.

Each spike is defined as a radial perturbation to the polygon at a
specified angular location $\theta_0$:
\begin{equation}
  \delta r(\theta) = h\,\exp\!\!\left(-\frac{(\theta-\theta_0)^2}{2\sigma_s^2}\right)
  \label{eq:spike_profile}
\end{equation}
where $h$ is the spike height (parameter to optimize),
$\sigma_s = w_{\mathrm{base}} / (2\sqrt{2\ln 2})$ is the Gaussian
width parameter, and $w_{\mathrm{base}}$ is the full-width-at-half-maximum
angular extent.  A narrow angular width $w_{\mathrm{base}} = 10^\circ$
is used to maximize field concentration effect per unit area added to
the polygon.  Spikes are placed at $\theta_0 = 0^\circ$ and
$180^\circ$ (the two parallel-to-$\vect{E}$ face midpoints of the
square).

\subsection{Turntable Rotation Model}
\label{sec:turntable_method}

The physical scenario is a motorized turntable that rotates the entire
part assembly---the doped-region polygon rasterized on the powder
bed---at discrete angles within the fixed RF field (electrodes remain
stationary).  Rotation is discrete (not continuous) because a
continuous rotation would sweep each boundary element through all
field orientations simultaneously, which is equivalent to perfect
averaging and is the upper limit of what rotation can achieve.

At each rotation event at time $t_k$, the cumulative rotation angle
is incremented by $\Delta\theta = 90^\circ$, and the spatial fields
$(T, \rhorel, \phimelt)$ are remapped from the old to the new grid:
\begin{equation}
  f^{\mathrm{new}}(x, z) = f^{\mathrm{old}}(R_{-\Delta\theta}(x,z))
  \label{eq:rotation_remap}
\end{equation}
where $R_{-\Delta\theta}$ is the inverse rotation operator.  The
remapping uses bilinear interpolation via
\texttt{scipy.ndimage.map\_coordinates}, which preserves the spatial
structure of hot-spots and density gradients with sub-grid accuracy.

The rotation schedule is time-based rather than melt-fraction-based
(the latter would require real-time density sensing):
\begin{equation}
  t_k = \frac{k}{N_r + 1}\,t_{\mathrm{total}},
  \quad k = 1, 2, \ldots, N_r
  \label{eq:rotation_schedule}
\end{equation}
For $N_r = 4$ rotations over $t_{\mathrm{total}} = \SI{360}{\second}$:
rotations occur at $t = 72, 144, 216, 288\,\si{\second}$, each adding
$90^\circ$ for a total of $360^\circ$ over the full exposure.  This
schedule is uniformly spaced so that each of the four $90^\circ$
orientations occupies an equal $\SI{72}{\second}$ interval.

%% ---------------------------------------------------------------
\section{Canonical Baseline Results}
\label{sec:results_baseline}
%% ---------------------------------------------------------------

\subsection{EQS Field Distribution: 20 mm Square at 6 Minutes}

\Cref{fig:canonical_fields} presents the canonical simulation result
for a $20 \times \SI{20}{\milli\meter}$ square at $t = \SI{360}{\second}$.
The electric potential $|V|$ shows the expected near-linear gradient
between electrodes, slightly distorted by the high-permittivity doped
region.  The $\Qrf$ map reveals the characteristic non-uniformity
driven by the three mechanisms described in \cref{sec:physics}:
corners at $\approx 37\times$ the spatial mean, perpendicular faces at
$\approx 11\times$, and parallel faces at $\approx 1.6\times$.  The
numerical values of $\Qrf$ at representative locations are listed in
\cref{tab:qrf_values}.

\begin{figure}[tb]
  \centering
  \includegraphics[width=\columnwidth]{xz_uniform_6min_v2/electric_fields.png}
  \caption{Canonical 6-minute run: electric potential $|V|$, electric field
           magnitude $|\vect{E}|$, and volumetric RF heating rate $\Qrf$ for
           a $20\times\SI{20}{\milli\meter}$ square.  The $\approx 37\times$
           corner-to-center ratio in $\Qrf$ is evident.  Chamber:
           $60\times\SI{60}{\milli\meter}$; generator: \SI{500}{\watt}, 2\%
           coupling efficiency.}
  \label{fig:canonical_fields}
\end{figure}

\begin{table}[tb]
\centering
\caption{Volumetric RF heating rate $\Qrf$ at representative boundary
         locations for the 20~mm square at $t = \SI{360}{\second}$.
         Ratios are relative to the spatial mean over the doped region.}
\label{tab:qrf_values}
\begin{tabular}{lSS[table-format=2.0]}
\toprule
Location & {$\Qrf$ (\si{\watt\per\meter\cubed})} & {Ratio to mean} \\
\midrule
Corner & 2.26e9 & 37 \\
Top/bottom face midpoint & 3.34e8 & 11 \\
Left/right face midpoint & 4.77e7 & 1.6 \\
Part center & 3.13e7 & 1.0 \\
\bottomrule
\end{tabular}
\end{table}

\begin{figure*}[tb]
  \centering
  \includegraphics[width=\textwidth]{xz_uniform_6min_v2/paper_style_report.png}
  \caption{Five-panel canonical result for $20\times\SI{20}{\milli\meter}$
           square at $t = \SI{360}{\second}$ (6 min).  Left to right:
           temperature with E-field streamlines, volumetric heating rate
           $\Qrf$, melt fraction $\phimelt$, relative density $\rhorel$,
           and time-series evolution of mean temperature, melt fraction,
           and relative density.  The spatial non-uniformity in $\rhorel$
           ($\sigma_{\mathrm{bnd}} = 0.072$) driven by the corner/face
           $\Qrf$ gradient is clearly visible.}
  \label{fig:canonical_full}
\end{figure*}

\subsection{Thermal and Density Results}

\Cref{fig:canonical_full} presents the five-panel canonical result.
At $t = \SI{360}{\second}$:
\begin{itemize}
  \item Mean temperature: $\bar{T} = \SI{181.2}{\celsius}$;
        maximum: $T_{\max} = \SI{197.6}{\celsius}$
  \item Mean melt fraction: $\bar{\phimelt} = 0.626$;
        fraction of cells at $T \geq T_{\mathrm{pc}}$: 82.5\%
  \item Mean relative density: $\bar{\rhorel} = 0.610$
  \item Boundary density metrics: $\bndstd = 0.072$,
        $r_{\mathrm{bnd}} = 0.699$
\end{itemize}
The temperature field shows local hot-spots at corners reaching
$\SI{197.6}{\celsius}$ while parallel-face midpoints approach the
melt threshold from below.  The density field shows the same pattern
amplified by the non-linear densification kinetics: corners at
$\rhorel \approx 0.80$ vs. parallel-face midpoints at
$\rhorel \approx 0.55$.

\subsection{Time-Exposure Sweep: 6--10 Minutes}

\Cref{tab:time_sweep} presents the systematic time-exposure sweep,
demonstrating how all state variables evolve with exposure time for
the canonical square geometry.  The melt fraction $\bar{\phimelt}$
increases rapidly from 6 to 8 minutes (0.626 to 0.964) and saturates
near 1.0 by 9 minutes, while the relative density continues to increase
more slowly (0.610 to 0.932 over the full 4-minute window).  The
boundary density standard deviation $\bndstd$ monotonically increases
with exposure time: as the already-hot corners advance into the
liquid-state sintering regime and the parallel faces merely enter
solid-state sintering, the density gap between these regions widens.

\begin{table}[tb]
\centering
\caption{Time-exposure sweep for the canonical $20\times\SI{20}{\milli\meter}$
         square at \SI{500}{\watt}, 2\% coupling efficiency.  $\bar{\phimelt}$
         is the spatial mean melt fraction over the part domain;
         $\bar{\rhorel}$ is the spatial mean relative density.}
\label{tab:time_sweep}
\begin{tabular}{cSSSSS}
\toprule
$t$ (min) & {$\bar{T}$ (\si{\celsius})} & {$T_{\max}$ (\si{\celsius})}
           & {$\bar{\phimelt}$} & {$\bar{\rhorel}$} & {$\bndstd$} \\
\midrule
6  & 181.2 & 197.6 & 0.626 & 0.610 & 0.0724 \\
7  & 185.9 & 205.9 & 0.894 & 0.692 & 0.0763 \\
8  & 195.2 & 214.1 & 0.964 & 0.783 & 0.0801 \\
9  & 205.0 & 222.6 & 0.996 & 0.868 & 0.0822 \\
10 & 214.9 & 231.3 & 1.000 & 0.932 & 0.0807 \\
\bottomrule
\end{tabular}
\end{table}

At 10 minutes, $\bndstd$ decreases slightly from its peak at 9 minutes.
This reflects the increasing role of thermal conduction: once corners
are fully melted, additional heat flows conductively to face midpoints,
partially closing the density gap.  However, even at 10 minutes,
$\bndstd = 0.081$ remains 12\% higher than at 6 minutes, confirming
that extended exposure worsens rather than ameliorates density
non-uniformity for the geometries and conditions studied.

%% ---------------------------------------------------------------
\section{Primitive Shape Library}
\label{sec:results_shapes}
%% ---------------------------------------------------------------

\Cref{fig:primitive_sweep} presents the results of simulating 17
primitive shapes under identical conditions (6 minutes, \SI{500}{\watt},
same chamber and grid).  \Cref{tab:shape_library} summarizes the
key metrics.

\begin{figure*}[tb]
  \centering
  \includegraphics[width=\textwidth]{primitive_sweep.png}
  \caption{Primitive shape library: $\rhorel$ maps at $t = \SI{360}{\second}$
           for 17 shapes under identical conditions (\SI{500}{\watt},
           2\% coupling efficiency, 6-minute exposure).  All shapes are
           approximately $20\times\SI{20}{\milli\meter}$ bounding box.
           The strong heating anisotropy (corner/vertex overheating,
           parallel-face underheating) is apparent across all shapes but
           manifests differently depending on shape symmetry and aspect
           ratio.}
  \label{fig:primitive_sweep}
\end{figure*}

\begin{table*}[tb]
\centering
\caption{Primitive shape simulation results at $t = \SI{360}{\second}$.
         $\bndstd$ is the boundary standard deviation (\cref{eq:bnd_std});
         $r_{\mathrm{bnd}}$ is the min/max boundary density ratio
         (\cref{eq:minmax_ratio}).  Shapes are sorted by $\bndstd$.}
\label{tab:shape_library}
\begin{tabular}{lcSSSSS}
\toprule
Shape & Vertices & {$\bndstd$} & {$r_{\mathrm{bnd}}$}
      & {$\bar{\rhorel}$} & {$\bar{T}$ (\si{\celsius})} & {$f_{\mathrm{melt}}$} \\
\midrule
Star-6       & 12  & 0.0551 & 0.832 & 0.975 & 258.6 & 0.95 \\
Octagon      & 8   & 0.0570 & 0.758 & 0.873 & 221.7 & 0.91 \\
Square       & 4   & 0.0724 & 0.699 & 0.610 & 181.2 & 0.83 \\
Rounded rect & 4*  & 0.1054 & 0.712 & 0.647 & 183.5 & 0.85 \\
Rectangle    & 4   & 0.1056 & 0.601 & 0.724 & 192.6 & 0.88 \\
Star-5       & 10  & 0.1437 & 0.550 & 0.958 & 266.5 & 0.96 \\
Trapezoid    & 4   & 0.1334 & 0.588 & 0.749 & 197.1 & 0.89 \\
Ellipse      & ---  & 0.1509 & 0.570 & 0.944 & 249.2 & 0.94 \\
Triangle     & 3   & 0.1231 & 0.602 & 0.884 & 227.8 & 0.92 \\
Equil.\ triangle & 3 & 0.1586 & 0.551 & 0.916 & 255.7 & 0.94 \\
L-shape      & 6   & 0.1597 & 0.575 & 0.771 & 198.4 & 0.88 \\
Hexagon      & 6   & 0.1798 & 0.550 & 0.853 & 217.9 & 0.91 \\
Pentagon     & 5   & 0.1795 & 0.550 & 0.870 & 224.5 & 0.92 \\
T-shape      & 8   & 0.1765 & 0.567 & 0.758 & 197.2 & 0.88 \\
Circle       & --- & 0.1661 & 0.558 & 0.773 & 202.2 & 0.88 \\
Diamond      & 4   & 0.1689 & 0.550 & 0.644 & 179.5 & 0.81 \\
Cross        & 12  & 0.1908 & 0.554 & 0.812 & 210.2 & 0.89 \\
\bottomrule
\multicolumn{7}{l}{\small *Rounded rectangle vertices count is nominal; the shape has
                   continuous curvature at corners.}
\end{tabular}
\end{table*}

\subsection{Effect of Shape Symmetry}

The Star-6 and Octagon achieve the lowest $\bndstd$ (0.055 and 0.057,
respectively) of all 17 shapes, despite operating at far higher mean
temperatures than the square.  This reflects their higher degree of
rotational symmetry: both shapes present many faces at roughly equal
angles to the applied field, distributing E-field concentration more
evenly.  The Star-6's 12 vertices each act as local field concentrators,
but because they are uniformly distributed in angle, no single region
is preferentially over- or under-heated.

In contrast, the Cross, T-shape, and L-shape achieve the highest
$\bndstd$ (0.191, 0.177, 0.160, respectively).  These shapes have
multiple extended parallel-to-$\vect{E}$ face segments separated by
re-entrant concave corners that block field-line access.  The
combination of large under-heated face area and limited heat diffusion
from nearby hot regions creates the largest density gradients.

\subsection{The Diamond Paradox}

The diamond (square rotated $45^\circ$) presents an instructive
case: it has the lowest mean temperature of any shape (only
$\SI{179.5}{\celsius}$, barely above the melt threshold) and the
lowest overall density ($\bar{\rhorel} = 0.644$), yet its tip corners
reach extreme temperatures ($T_{\max} \approx \SI{272}{\celsius}$).
The reason is geometric: a diamond oriented at $45^\circ$ presents
all four faces at $45^\circ$ to the applied field.  There is no
purely perpendicular face ($0^\circ$) and no purely parallel face
($90^\circ$); instead, all faces are at intermediate angles where the
effective field enhancement is $\approx 3$--$5\times$ the mean.  While
vertices still have the tip effect, the faces---which constitute the
majority of the boundary---receive field enhancement comparable to the
center rather than the strong normal-field enhancement of the $0^\circ$
faces.  The result is overall lower heating with extreme tip hot-spots
and poor overall densification.

\subsection{Shape Design Implications}

The shape library reveals three practical design guidelines for RFAM:
(a) shapes with many equal-angle vertices distribute field concentration
most uniformly;
(b) concave re-entrant features systematically increase $\bndstd$
regardless of overall shape symmetry (L, T, Cross);
(c) reducing corner sharpness (Rounded rect vs. Rectangle) reduces
$\bndstd$ from 0.106 to 0.105 and increases $r_{\mathrm{bnd}}$ from
0.601 to 0.712, even though the improvement in $\bndstd$ itself is
marginal---the primary benefit is the increase in the worst-case
minimum boundary density.

%% ---------------------------------------------------------------
\section{Geometry Prewarp Results}
\label{sec:results_prewarp}
%% ---------------------------------------------------------------

\subsection{$\Leff$ Calibration for Circle and Square}

\Cref{tab:leff_calibration} presents the calibration sweep results
for the circle and square.  For the circle, the optimal $\Leff =
\SI{5}{\milli\meter}$ achieves $\RMSE = \SI{0.103}{\milli\meter}$ in
2 iterations.  The circle's smooth boundary with gradual EPE variation
enables the algorithm to converge reliably to a low RMSE; large
$\Leff$ values ($\geq \SI{10}{\milli\meter}$) lead to over-correction
oscillations that prevent convergence within the iteration budget.

For the square, the optimal $\Leff = \SI{7}{\milli\meter}$ achieves
$\RMSE = \SI{0.153}{\milli\meter}$.  The higher optimal $\Leff$ and
residual RMSE compared to the circle reflects the square's sharp corners,
where the EPE changes discontinuously at the $90^\circ$ vertex.  The
prewarp algorithm cannot fully correct these corner singularities
because the shrinkage model (\cref{eq:shrinkage_ratio}) does not capture
the three-dimensional redistribution of material at a sharp corner during
sintering.

\begin{table}[tb]
\centering
\caption{$\Leff$ calibration sweep: post-convergence geometric RMSE
         (\si{\milli\meter}) for circle and square targets.  Bold
         indicates the optimal value selected for each shape.}
\label{tab:leff_calibration}
\begin{tabular}{cSS}
\toprule
$\Leff$ (mm) & {Circle RMSE (mm)} & {Square RMSE (mm)} \\
\midrule
3  & 0.117 & 0.209 \\
\textbf{5}  & \textbf{0.103} & 0.182 \\
7  & 0.194 & \textbf{0.153} \\
10 & 0.167 & 0.174 \\
13 & 0.151 & 0.227 \\
16 & 0.151 & 0.241 \\
\bottomrule
\end{tabular}
\end{table}

\begin{figure}[tb]
  \centering
  \begin{subfigure}[b]{\columnwidth}
    \includegraphics[width=\columnwidth]{prewarp_calibrate_circle/calibration_sweep.png}
    \caption{Circle: $\RMSE$ vs.\ $\Leff$ at convergence.}
    \label{fig:leff_circle}
  \end{subfigure}
  \par\smallskip
  \begin{subfigure}[b]{\columnwidth}
    \includegraphics[width=\columnwidth]{prewarp_calibrate_square/calibration_sweep.png}
    \caption{Square: $\RMSE$ vs.\ $\Leff$ at convergence.}
    \label{fig:leff_square}
  \end{subfigure}
  \caption{Effective shrinkage length $\Leff$ calibration sweeps.
           Each point represents the post-convergence RMSE after running
           the full EPE-gradient-descent algorithm.  The optimal values
           are $\Leff = \SI{5}{\milli\meter}$ (circle) and
           $\Leff = \SI{7}{\milli\meter}$ (square).}
  \label{fig:leff_calibration}
\end{figure}

\subsection{Square Prewarp Validation}

\Cref{fig:square_prewarp} shows the prewarp result for the square
target.  The prewarp polygon (input to the simulator) is noticeably
convex: the corners protrude inward relative to a square, and the face
midpoints protrude outward slightly, compensating for the greater
shrinkage at corners vs. face midpoints.  The RMSE improves from
\SI{0.37}{\milli\meter} (uncorrected, iteration 0) to
\SI{0.153}{\milli\meter} (converged, iteration 2), a 59\% reduction.

The boundary density after prewarp is also improved: the prewarp polygon
forces the parallel-face midpoints to extend further outward, increasing
the local field coupling at those previously under-heated regions.  This
geometric effect reduces $\bndstd$ from 0.072 (unprewarped) to 0.040
(prewarped), a 44\% reduction---a genuine side benefit of the prewarp
in addition to its intended geometric correction.

\begin{figure}[tb]
  \centering
  \includegraphics[width=\columnwidth]{prewarp_square_v2/prewarp_report.png}
  \caption{Square prewarp result using $\Leff = \SI{7}{\milli\meter}$.
           Top: prewarp polygon (blue) vs.\ target (dashed red).
           Bottom: predicted sintered boundary (green) vs.\ target.
           RMSE reduces from \SI{0.37}{\milli\meter} (unprewarped) to
           \SI{0.153}{\milli\meter} after 2 iterations.}
  \label{fig:square_prewarp}
\end{figure}

\subsection{H-Shape Prewarp}
\label{sec:hshape_prewarp}

The H-shape presents a substantially more challenging prewarp problem.
Its bounding box is $20 \times \SI{20}{\milli\meter}$ with two vertical
legs (\SI{6}{\milli\meter} wide) connected by a horizontal crossbar
(\SI{6}{\milli\meter} tall, spanning between the legs at mid-height).
The shape has 5 geometrically distinct face categories:
\begin{enumerate}
  \item Outer vertical legs (parallel to $\vect{E}$, under-heated):
        $\rhorel^{\mathrm{bnd}} \approx 0.610$ at midpoints, as low
        as 0.552 at worst cells
  \item Outer top/bottom faces (perpendicular to $\vect{E}$,
        better-heated): $\rhorel^{\mathrm{bnd}} \approx 0.729$
  \item Crossbar top/bottom faces (perpendicular, partially shielded
        by the legs): $\rhorel^{\mathrm{bnd}} \approx 0.632$
  \item Inner notch faces (parallel, further shielded): $\rhorel \approx 0.679$
  \item Concave inside corners (locally over-heated but inaccessible
        to field from outside due to the notch geometry)
\end{enumerate}

Without correction, $\RMSE = \SI{0.682}{\milli\meter}$ and
$\bndstd = 0.078$.  After 2 prewarp iterations with
$\Leff = \SI{7}{\milli\meter}$, the RMSE decreases to
\SI{0.221}{\milli\meter} (68\% reduction) and $\bndstd$ decreases
to 0.056 (28\% reduction).  The prewarp polygon expands to a bounding
box of $21.00 \times \SI{20.90}{\milli\meter}$, approximately 5\%
larger than the target in each dimension.

\Cref{fig:hshape_prewarp} shows the H-shape prewarp result.  The
non-symmetric expansion is clearly visible: the outer vertical leg
faces protrude more than the horizontal faces, correctly anticipating
the greater shrinkage at those under-heated locations.

\begin{figure}[tb]
  \centering
  \includegraphics[width=\columnwidth]{prewarp_H/prewarp_report.png}
  \caption{H-shape prewarp result ($\Leff = \SI{7}{\milli\meter}$).
           The prewarp polygon (blue) is asymmetrically expanded relative
           to the target (dashed red), with larger expansion at the
           parallel-to-$\vect{E}$ vertical leg faces that would otherwise
           under-sinter and shrink less than expected.  RMSE: 0.682
           (no correction) $\to$ \SI{0.221}{\milli\meter} after
           2 iterations.}
  \label{fig:hshape_prewarp}
\end{figure}

%% ---------------------------------------------------------------
\section{Spike Assist Features}
\label{sec:results_spikes}
%% ---------------------------------------------------------------

\Cref{tab:spike_sweep} presents the spike height sweep results for
the prewarped square polygon.  The spike profile (\cref{eq:spike_profile})
with $w_{\mathrm{base}} = 10^\circ$ is placed at the two
parallel-to-$\vect{E}$ face midpoints ($\theta_0 = 0^\circ$ and
$180^\circ$).

\begin{table}[tb]
\centering
\caption{Spike assist feature height sweep on the prewarped square
         polygon.  Spikes placed at $\theta_0 = 0^\circ$ and $180^\circ$,
         $w_{\mathrm{base}} = 10^\circ$.  Baseline (h=0) is the
         prewarped square without spikes.}
\label{tab:spike_sweep}
\begin{tabular}{cSSS}
\toprule
$h$ (mm) & {$\bndstd$} & {$r_{\mathrm{bnd}}$} & {Improvement (\%)} \\
\midrule
0.0 & 0.0403 & 0.799 & 0.0 \\
1.6 & 0.0405 & 0.796 & -0.4 \\
3.2 & 0.0388 & 0.798 & +3.8 \\
4.8 & 0.0379 & 0.797 & +5.9 \\
6.4 & 0.0365 & 0.801 & +9.5 \\
7.6 & 0.0358 & 0.802 & +11.0 \\
7.8 & 0.0354 & 0.802 & +12.1 \\
7.9 & 0.0349 & 0.802 & +13.5 \\
\textbf{8.0} & \textbf{0.0347} & \textbf{0.802} & \textbf{+14.1} \\
\bottomrule
\end{tabular}
\end{table}

\begin{figure}[tb]
  \centering
  \includegraphics[width=\columnwidth]{spike_coopt_v2_square/coopt_report.png}
  \caption{Spike assist feature optimization for the square.  Top:
           optimized spiked polygon ($h = \SI{8}{\milli\meter}$)
           overlaid on the target.  Bottom: $\rhorel$ field showing
           modest improvement at the parallel-face midpoints.  The
           $\bndstd$ improvement is only 14.1\% at the maximum practical
           spike height.}
  \label{fig:spike_coopt}
\end{figure}

Spike features improve $\bndstd$ monotonically with height, reaching
14.1\% improvement at $h = \SI{8}{\milli\meter}$.  However, three
limitations prevent spikes from being practically useful:

\paragraph{Limited spatial reach.}
Each $10^\circ$ spike covers approximately 2 boundary cells on a
$120\times120$ grid.  The 160-cell square perimeter therefore has
only $\approx 2.5\%$ of its boundary cells directly affected by each
spike.  The concentration effect at the spike tip does not diffuse
meaningfully to neighboring face cells within the $\SI{360}{\second}$
exposure time.

\paragraph{Geometric impracticality.}
At $h = \SI{8}{\milli\meter}$ (40\% of the $\SI{20}{\milli\meter}$ part
dimension), the spike would clearly print as a real geometric feature.
An SRAF in optical lithography is deliberately below the resist
resolution threshold; no such threshold exists for RFAM, where any
conductive feature will sinter.  A spike of this height would appear
as a $\SI{8}{\milli\meter}$-tall antenna protrusion on the finished
part---unacceptable for any functional application.

\paragraph{Geometric RMSE penalty.}
At $h = \SI{8}{\milli\meter}$, the geometric RMSE of the spike-modified
polygon relative to the square target is $\approx \SI{2.1}{\milli\meter}$,
compared to \SI{0.153}{\milli\meter} for the prewarp alone.  The spike
thus sacrifices the primary benefit of the prewarp algorithm in exchange
for a modest secondary benefit.

The conclusion is that spike assist features are not a viable
stand-alone strategy for RFAM density uniformity improvement.  Their
utility is limited by the fundamental difference between optical
lithography (where resolution limits enforce feature size constraints)
and RFAM (where all sintering is above the resolution threshold).

%% ---------------------------------------------------------------
\section{Turntable Rotation Strategy}
\label{sec:results_turntable}
%% ---------------------------------------------------------------

\subsection{Physical Mechanism}

Rotating the part $90^\circ$ converts the two parallel-to-$\vect{E}$
faces into perpendicular-to-$\vect{E}$ faces and vice versa.  With
four $90^\circ$ rotations, each face spends equal time in both
orientations.  In the limit of many infinitesimal rotations, the
exposure would be perfectly uniform (azimuthally averaged field), but
any finite number of discrete rotations retains some residual anisotropy.

The time-based schedule (\cref{eq:rotation_schedule}) distributes
rotations uniformly throughout the exposure.  A single rotation at
$t = \SI{180}{\second}$ (midpoint) would allow the first three minutes
of exposure at one orientation to dominate, while the second three
minutes at the rotated orientation merely attenuate the existing
non-uniformity partially.  By contrast, four rotations at 72 s
intervals ensure that no single orientation dominates: each of the
four $90^\circ$ orientations is active for exactly $\SI{72}{\second}$,
\textit{and} each rotation occurs before significant melting has
occurred (since $T$ only approaches $\Tpc$ near the end of the
exposure for the 6-minute schedule).

\subsection{Results}

\Cref{tab:turntable_comparison} presents the comprehensive comparison
of rotation strategies applied to the square.

\begin{table}[tb]
\centering
\caption{Boundary density uniformity metrics for the square under
         various rotation strategies.  \% improvement is relative to
         the baseline (no rotation, no prewarp).}
\label{tab:turntable_comparison}
\begin{tabular}{lcSS}
\toprule
Strategy & Prewarp & {$\bndstd$} & {$r_{\mathrm{bnd}}$} \\
\midrule
No rotation (baseline) & No  & 0.0724 & 0.699 \\
1$\times$ 90° at $t=3$ min & No  & 0.0705 & 0.698 \\
4$\times$ 90° (time-based) & No  & 0.0585 & 0.700 \\
Prewarp only            & Yes & 0.0403 & 0.799 \\
Prewarp + 4$\times$ 90° & Yes & \textbf{0.0253} & \textbf{0.865} \\
\bottomrule
\end{tabular}
\end{table}

The single rotation at $t = \SI{3}{\minute}$ provides only 2.6\%
improvement in $\bndstd$ because the part has already completed 90\%
of its heating at the original orientation by the time the rotation
occurs.  In contrast, the 4-rotation schedule (rotations at 72, 144,
216, 288 s) achieves 19.2\% improvement, confirming the importance of
early, frequent rotation before significant melting.

The prewarp-only strategy reduces $\bndstd$ by 44\%, substantially
outperforming the turntable alone (19\%).  Combining prewarp with
4-rotation turntable achieves 65\% improvement ($\bndstd = 0.025$,
$r_{\mathrm{bnd}} = 0.865$).

\begin{figure*}[tb]
  \centering
  \includegraphics[width=\textwidth]{uniformity_comparison.png}
  \caption{Boundary density uniformity comparison for the square:
           five strategies from baseline (no correction) to combined
           prewarp + 4$\times$ turntable.  Bar charts show $\bndstd$
           and $r_{\mathrm{bnd}}$ for each strategy; spatial $\rhorel$
           maps illustrate the progressively more uniform boundary
           densification.  The combined strategy achieves $\bndstd =
           0.025$ and $r_{\mathrm{bnd}} = 0.865$, representing 65\%
           improvement over the baseline.}
  \label{fig:uniformity_comparison}
\end{figure*}

\Cref{fig:uniformity_comparison} provides a visual comparison of all
strategies.  The spatial $\rhorel$ maps confirm that the combined
approach produces a near-uniform distribution across the entire
boundary of the square, with the worst-sintered boundary cells
($\rhorel \approx 0.55$ at the parallel-face midpoints in the
baseline) improving to $\rhorel \approx 0.62$ in the combined case.

%% ---------------------------------------------------------------
\section{H-Shape: Combined Prewarp and Turntable}
\label{sec:results_combined}
%% ---------------------------------------------------------------

\subsection{H-Shape Baseline and Prewarp}

The H-shape is selected as the primary complex validation case for the
following reasons: (1) it is topologically non-trivial with an internal
cutout; (2) it has 5 geometrically distinct face types with different
electromagnetic coupling; and (3) it is non-symmetric with respect to
the $\vect{E}$-field direction (one axis of symmetry parallel to
$\vect{E}$, one perpendicular).

\Cref{tab:hshape_comparison} presents the progressive improvement
from baseline through combined correction.  Without any correction,
the H-shape has $\bndstd = 0.078$ and $r_{\mathrm{bnd}} = 0.684$,
representing a boundary density spread from 0.552 (worst parallel-face
midpoint) to 0.806 (best corner cell), a 46\% span.

\begin{table}[tb]
\centering
\caption{H-shape boundary density metrics under progressive correction
         strategies.  Improvement is relative to the H-shape baseline.}
\label{tab:hshape_comparison}
\begin{tabular}{lcSS}
\toprule
Strategy & Prewarp & {$\bndstd$} & {$r_{\mathrm{bnd}}$} \\
\midrule
Baseline (no correction)      & No  & 0.0781 & 0.684 \\
Prewarp only ($\Leff = 7$ mm) & Yes & 0.0562 & 0.742 \\
4$\times$ 90° turntable only  & No  & 0.0516 & 0.721 \\
Prewarp + 4$\times$ 90°       & Yes & \textbf{0.0108} & \textbf{0.938} \\
\bottomrule
\end{tabular}
\end{table}

\begin{figure*}[tb]
  \centering
  \includegraphics[width=\textwidth]{H_uniformity_comparison.png}
  \caption{H-shape boundary density uniformity comparison.  Left to
           right: baseline, prewarp only, turntable only, combined
           prewarp + 4$\times$ turntable.  The combined strategy achieves
           $\bndstd = 0.011$ and $r_{\mathrm{bnd}} = 0.938$: an 86\%
           reduction in boundary density non-uniformity.  Boundary cells
           are uniformly sintered to within $\pm 1.8\%$ of the mean
           boundary density.}
  \label{fig:hshape_comparison}
\end{figure*}

\begin{figure*}[tb]
  \centering
  \includegraphics[width=\textwidth]{prewarp_H_turntable_4x_eval/paper_style_report.png}
  \caption{Combined prewarp + 4$\times$ turntable for the H-shape:
           temperature, $\Qrf$, melt fraction, and relative density fields
           at $t = \SI{360}{\second}$.  The rotational averaging
           is visible in the near-symmetric $\rhorel$ distribution,
           where all boundary regions achieve comparable densification
           despite the inherent anisotropy of the RF field.}
  \label{fig:hshape_combined}
\end{figure*}

\subsection{Combined H-Shape Results}

The combined prewarp + 4$\times$ turntable result for the H-shape
achieves:
\begin{itemize}
  \item $\bndstd = 0.0108$ (86.2\% reduction vs.\ baseline)
  \item $r_{\mathrm{bnd}} = 0.938$ (vs.\ 0.684 baseline)
  \item Boundary density range: $[0.550, 0.586]$ --- all boundary cells
        within 3.6\% of each other
  \item Mean temperature: $\bar{T} = \SI{179.9}{\celsius}$
\end{itemize}
\Cref{fig:hshape_comparison,fig:hshape_combined} illustrate the spatial
density distribution and the 5-panel field summary.

\subsection{Synergy Analysis: Why Combined $>$ Prewarp $+$ Turntable}

If prewarp and turntable were independent mechanisms (i.e., their
effects did not interact), the combined improvement would be
approximately the sum of their individual improvements.  For the
H-shape:
\begin{itemize}
  \item Prewarp alone: 28\% reduction in $\bndstd$
  \item Turntable alone: 34\% reduction in $\bndstd$
  \item Naive sum: 62\% expected
  \item Actual combined: 86\%
\end{itemize}
The 24-percentage-point surplus indicates a positive synergy.

The synergy arises from the following structural interaction:
\begin{enumerate}
  \item The prewarp polygon is designed for the field-anisotropic
        baseline configuration---it extends the parallel-face regions
        further outward to compensate for their lower shrinkage.  This
        creates a modified polygon with more surface area on those faces.
  \item The additional surface area at the parallel-face regions
        improves their field coupling: a face that protrudes further
        from the center encounters a slightly higher $|\vect{E}|$ and
        presents more area to field-line flux, increasing local $\Qrf$.
  \item The turntable then rotates this already-optimized polygon through
        all four $90^\circ$ orientations.  At each orientation, the
        extended face regions---which were designed for the worst-case
        (parallel) orientation---benefit equally when they rotate to
        the perpendicular orientation where they receive $11\times$
        mean coupling.
  \item The net result is that every face region alternates between
        one interval of modest coupling (parallel) and one interval of
        strong coupling (perpendicular), with the prewarp ensuring that
        the total shrinkage in each case lands exactly at the target.
\end{enumerate}
In contrast, without prewarp, the turntable rotates the unmodified
H-shape, which has equal-length faces in all orientations.  The
turntable helps, but without the geometric optimization, no face region
has a structural advantage in the parallel orientation.

\subsection{Physical Limits of the Combined Approach}

The $r_{\mathrm{bnd}} = 0.938$ achieved in the combined case implies
that boundary density ranges over only $[0.550, 0.586]$---a 6.5\% span.
The residual non-uniformity comes from two sources:

\paragraph{Discrete rotation residual.}
With 4 rotations (5 intervals of $\SI{72}{\second}$ each), one
orientation is sampled at the $0^\circ$ starting position and the
four rotated positions.  The H-shape has twofold rotational symmetry
(not fourfold), so a $90^\circ$ rotation produces a distinct field
pattern.  The resulting exposure is not perfectly uniform; a small
residual anisotropy survives.

\paragraph{Concave corner shielding.}
The H-shape's internal notch corners are shielded from the external
field regardless of rotation angle.  These cells receive heating only
by conduction from neighboring cells, which is insufficient to fully
densify them in 6 minutes.  Increasing exposure time would improve
these cells at the cost of overheating the outer corners.

%% ---------------------------------------------------------------
\section{Discussion}
\label{sec:discussion}
%% ---------------------------------------------------------------

\subsection{Model Validation and Calibration}

The forward model requires only two calibrated parameters: $\sigma =
\SI{0.04}{\siemens\per\meter}$ and $\eta = 0.020$.  Both are
physically motivated: the doped-powder conductivity is a material
property that can in principle be measured independently, and the
coupling efficiency $\eta$ is a property of the electrode geometry and
impedance matching network.  The agreement with COMSOL (which uses a
spatially varying $\sigma$ obtained from a 3D interpolation function)
in both the spatial pattern of $\Qrf$ and the steady-state temperature
validates the model for the purposes of the prewarp and uniformity
optimization studies presented here.

The key simplification relative to COMSOL is the uniform $\sigma$
approximation.  The COMSOL Tuned\_Sigma.mph model uses a spatially
varying conductivity with the characteristic \emph{shell} pattern
(high at outer faces/edges, low at the center), resulting in an
average $\sigma \ll \SI{0.04}{\siemens\per\meter}$.  The Python model
compensates for this by setting $\eta = 0.020$, which reduces the
effective absorbed power to match the COMSOL result.  The spatial
pattern of $\Qrf$ in both models shows the same corner-to-face-to-center
hierarchy, confirming that the dominant spatial structure is driven by
geometry (boundary conditions) rather than by the $\sigma$ distribution.

\subsection{Comparison to Optical ILT}

The EPE-based prewarp algorithm is directly analogous to standard ILT
methods~\cite{Jia2014ilt,Shen2009ilt}, with three key differences:

\paragraph{Forward model complexity.}
Optical ILT uses an analytical or semi-analytical forward model for the
aerial image (Hopkins or Abbe formulation); this enables analytic
gradient computation and fast optimization.  The RFAM forward model
requires a full PDE solve ($\approx$\SI{30}{\second} per iteration on a
laptop).  Consequently, the RFAM prewarp can afford only 2--3 iterations,
whereas optical ILT may run hundreds.

\paragraph{Binary vs.\ continuous mask.}
Optical masks are binary (opaque/transparent); the prewarp polygon is
a continuous boundary.  The RFAM analog of the mask is the input
polygon, which is free to take any shape---there is no tone
constraint.  This makes the RFAM problem in some ways easier
(continuous design space) and in some ways harder (no natural
binarization to prevent over-complicated geometries).

\paragraph{Shrinkage nonlinearity.}
The optical aerial image is related to the mask by a linear
convolution; the sintered boundary is related to the input polygon by
a nonlinear, history-dependent PDE system.  The Arrhenius densification
kinetics and phase-change latent heat introduce strong nonlinearities
that prevent the use of gradient-based optical ILT algorithms without
significant modification.

\subsection{Practical Implementation Considerations}

\paragraph{Rotation actuator.}
A simple stepper-motor-driven turntable below the powder bed would
implement the 4-rotation schedule.  The total rotation ($360^\circ$
over $\SI{360}{\second}$) is not demanding: $1^\circ$/s average.  The
rotation need not be continuous; a discrete 4-step schedule with
72-second holds is straightforward to implement with a low-cost
stepper motor and controller.

\paragraph{Prewarp pipeline.}
In a production RFAM workflow, the prewarp algorithm would run
offline prior to printing.  The 2--3 forward model evaluations required
for convergence ($\approx\SI{1}{\minute}$ each on the target hardware)
represent a modest pre-computation cost.  The converged prewarp polygon
replaces the original part outline in the binder-jetting layer mask,
requiring no changes to the RF exposure hardware.

\paragraph{3D extension.}
The current model is 2D (cross-section in the XZ plane).  Real parts
have a finite out-of-plane dimension and 3D E-field structure.  The
2D model is exact for infinitely long parts in the Y-direction and is
an approximation for parts with finite aspect ratios.  The 3D
generalization would require a 3D EQS solve (much higher computational
cost) and a 3D prewarp algorithm (vertex update in 3D with surface
normal computation on a mesh).  The turntable strategy extends directly
to 3D by rotating about the Z-axis (perpendicular to the electrode
faces).

%% ---------------------------------------------------------------
\section{Conclusions}
\label{sec:conclusions}
%% ---------------------------------------------------------------

This paper has presented and evaluated three strategies for improving
dimensional accuracy and sintering uniformity in Radio Frequency
Additive Manufacturing.  The following conclusions are supported by
the simulation results:

\begin{enumerate}

\item \textbf{EQS forward model accuracy.}  A 2D EQS+thermal+densification
      model with uniform $\sigma = \SI{0.04}{\siemens\per\meter}$ and
      $2\%$ coupling efficiency reproduces the physically expected
      corner/face/center heating hierarchy ($37\times$ / $11\times$ /
      $1.6\times$ the mean) derived from electromagnetic boundary
      conditions.  Only two calibrated parameters are required, both
      with physical interpretations.

\item \textbf{EPE prewarp convergence.}  The EPE-gradient-descent
      prewarp algorithm converges in 2--3 iterations to geometric RMSE
      below \SI{0.25}{\milli\meter} for simple shapes (circle, square)
      and below \SI{0.30}{\milli\meter} for complex shapes (H), starting
      from geometric RMSE of 0.37--\SI{0.68}{\milli\meter}.  This
      represents a 59--68\% reduction in geometric RMSE.

\item \textbf{$\Leff$ calibration.}  The optimal effective shrinkage
      length is $\SI{5}{\milli\meter}$ for the circle and
      $\SI{7}{\milli\meter}$ for the square and H-shape.  Physically,
      $\Leff$ relates to the characteristic sintering diffusion length
      at 6 minutes exposure; the higher value for shapes with sharp
      corners reflects the three-dimensional material redistribution
      at vertices that is not captured by the 2D model.

\item \textbf{Shape library insights.}  High-symmetry shapes (Star-6,
      Octagon) achieve the lowest boundary density non-uniformity
      ($\bndstd \approx 0.055$) by distributing field concentration
      uniformly.  Concave shapes (Cross, T, L) have the highest
      $\bndstd$ ($\approx 0.19$).  These results provide quantitative
      design guidance for RFAM part geometry selection.

\item \textbf{Spike assist features are limited.}  The RFAM analog of
      optical SRAFs---narrow convex protrusions at under-heated face
      midpoints---provides at most 14\% improvement in $\bndstd$ at
      geometrically impractical spike heights.  Unlike optical SRAFs,
      RFAM spikes have no resolution threshold below which they would
      not print, making them structurally incompatible with the SRAF
      concept as applied in optics.

\item \textbf{Turntable rotation is effective.}  Four $90^\circ$
      rotations at equal time intervals reduces $\bndstd$ by 19\% for
      the bare square and 34\% for the H-shape.  A single late rotation
      is nearly useless (2.6\% improvement), confirming the necessity
      of the early, uniform schedule.

\item \textbf{Synergistic combination achieves 86\% improvement.}
      Combining prewarp with 4-rotation turntable for the H-shape
      reduces $\bndstd$ from 0.078 to 0.011 (86.2\%) and increases
      $r_{\mathrm{bnd}}$ from 0.684 to 0.938.  The combined improvement
      exceeds the sum of independent improvements by 24 percentage
      points, owing to a structural interaction between the prewarp
      polygon and the turntable's exposure averaging.  All H-shape
      boundary cells sinter to within 3.6\% of each other.

\end{enumerate}

Future work will address: (a) 3D extension of the EQS and prewarp
models; (b) experimental validation against sintered Nylon~12 parts;
(c) parameter identification of the Arrhenius densification constants
from dilatometry measurements; and (d) continuous rotation strategies
and their theoretical limit of uniform exposure.

%% ---------------------------------------------------------------
\section*{Acknowledgments}
%% ---------------------------------------------------------------

The author thanks Prof.\ Christopher Seepersad for guidance on RFAM
process physics and C.B.\ Allison for the COMSOL reference model.
This work was supported by the Georgia Institute of Technology School
of Mechanical Engineering.

%% ---------------------------------------------------------------
%% References
%% ---------------------------------------------------------------
\begin{thebibliography}{99}

\bibitem{allison2015rfam}
C.B.~Allison,
``Radio frequency additive manufacturing,''
Ph.D.\ dissertation, The University of Texas at Austin, Austin, TX,
2015.

\bibitem{seepersad2014volumetric}
C.C.~Seepersad, C.B.~Allison, and J.J.~Beaman,
``Volumetric energy deposition in powder-bed additive manufacturing
using radio-frequency electric fields,''
in \textit{Proc.\ Solid Freeform Fabrication Symp.}, Austin, TX, 2014,
pp.~1--12.

\bibitem{seepersad2019rfam}
C.C.~Seepersad, T.~Mun, and R.~Lakshmi,
``Radio frequency additive manufacturing: selective sintering of
Nylon-12 using 27 MHz ISM-band energy,''
in \textit{Proc.\ Solid Freeform Fabrication Symp.}, Austin, TX, 2019,
pp.~1832--1843.

\bibitem{mccoy2025rfam_sim}
M.~McCoy,
``Two-dimensional electro-quasi-static and thermal simulation of
radio frequency additive manufacturing,''
Technical Report, Georgia Institute of Technology, 2025.

\bibitem{Wong2001ilt}
A.K.-K.~Wong,
\textit{Resolution Enhancement Techniques in Optical Lithography}.
Bellingham, WA: SPIE Press, 2001.

\bibitem{Poonawala2007ilt}
A.~Poonawala and P.~Milanfar,
``Mask design for optical microlithography---an inverse imaging
problem,''
\textit{IEEE Trans.\ Image Process.}, vol.~16, no.~3, pp.~774--788,
Mar.~2007.

\bibitem{Shen2009ilt}
Y.~Shen, N.~Jia, N.~Wong, and E.Y.~Lam,
``Robust level-set-based inverse lithography,''
\textit{Opt.\ Express}, vol.~19, no.~6, pp.~5511--5521, 2011.

\bibitem{Jia2014ilt}
N.~Jia, E.Y.~Lam, and M.~Smith,
``Fast pixel-based mask optimization for inverse lithography,''
\textit{J.\ Micro/Nanolithogr.\ MEMS MOEMS}, vol.~10, no.~4,
p.~043008, 2011.

\bibitem{Meixner1972singular}
J.~Meixner,
``The behavior of electromagnetic fields at edges,''
\textit{IEEE Trans.\ Antennas Propagat.}, vol.~20, no.~4,
pp.~442--446, Jul.~1972.

\bibitem{Frenkel1945}
J.~Frenkel,
``Viscous flow of crystalline bodies under the action of surface
tension,''
\textit{J.\ Phys.\ (Moscow)}, vol.~9, pp.~385--391, 1945.

\bibitem{Liebmann2001sraf}
L.~Liebmann, S.~Mansfield, A.~Wong, M.~Lavin, W.~Leipold, and
T.~Dunham,
``TCAD support for design rule development with lithography
simulation,''
in \textit{Proc.\ SPIE Optical Microlithography XIV},
vol.~4346, 2001, pp.~1--14.

\bibitem{Clark1996microwave}
D.E.~Clark and W.H.~Sutton,
``Microwave processing of materials,''
\textit{Annu.\ Rev.\ Mater.\ Sci.}, vol.~26, pp.~299--331, 1996.

\bibitem{Goodfellow2022nylon12}
Goodfellow Cambridge Ltd.,
``Nylon 12 (Polyamide 12) material data,''
Technical Datasheet, 2022.
[Online]. Available: \url{https://www.goodfellow.com/}

\bibitem{Laumer2014selective}
T.~Laumer, T.~Kroll, M.~Schmidt, and B.~Degner,
``Fundamental investigation of laser beam melting of polymers for
additive manufacture,''
\textit{J.\ Laser Appl.}, vol.~26, no.~4, p.~042003, 2014.

\bibitem{Drummer2010investigation}
D.~Drummer, D.~Rietzel, and F.~Kuehnlein,
``Development of a characterization approach for the sintering behavior
of new thermoplastics for selective laser sintering,''
\textit{Phys.\ Procedia}, vol.~5, pp.~533--542, 2010.

\end{thebibliography}

\end{document}
